\documentclass[11pt,a4paper]{article}

\usepackage[top=2.5cm, bottom=2.5cm, left=2.8cm, right=2.8cm]{geometry}
\usepackage{amsmath}
\usepackage{booktabs}
\usepackage{array}
\usepackage{xcolor}
\usepackage{hyperref}
\usepackage{microtype}
\usepackage{enumitem}
\usepackage{titlesec}
\usepackage{parskip}
\usepackage{fancyhdr}
\usepackage{lmodern}
\usepackage[T1]{fontenc}

% ── Colours ──────────────────────────────────────────────────────────────────
\definecolor{protoblue}{HTML}{1a4f7a}
\definecolor{rulecolor}{HTML}{1a4f7a}

% ── Section formatting ────────────────────────────────────────────────────────
\titleformat{\section}
  {\large\bfseries\color{protoblue}}
  {\thesection}{1em}{}[\color{rulecolor}\titlerule]
\titleformat{\subsection}
  {\normalsize\bfseries\color{protoblue}}
  {\thesubsection}{1em}{}
\titlespacing{\section}{0pt}{1.4em}{0.6em}
\titlespacing{\subsection}{0pt}{1em}{0.4em}

% ── Lists ─────────────────────────────────────────────────────────────────────
\setlist[itemize]{itemsep=2pt,topsep=4pt,leftmargin=1.4em}
\setlist[enumerate]{itemsep=2pt,topsep=4pt,leftmargin=1.6em}

% ── Header / footer ──────────────────────────────────────────────────────────
\pagestyle{fancy}
\fancyhf{}
\fancyhead[L]{\small\color{gray}N-Party Cyclic Atomic Swap --- Overview}
\fancyhead[R]{\small\color{gray}\thepage}
\renewcommand{\headrulewidth}{0.3pt}

% ── Hyperlinks ────────────────────────────────────────────────────────────────
\hypersetup{
  colorlinks=true,
  linkcolor=protoblue,
  urlcolor=protoblue,
  pdftitle={N-Party Cyclic Atomic Swap --- Overview and Assumptions}
}

% ─────────────────────────────────────────────────────────────────────────────
\begin{document}

% ── Title ────────────────────────────────────────────────────────────────────
\begin{titlepage}
\centering
\vspace*{2.5cm}
{\huge\bfseries\color{protoblue}
  N-Party Cyclic Atomic Swap\par}
\vspace{0.6cm}
{\Large\color{black}
  Overview and Assumptions\par}
\vspace{0.4cm}
{\large\color{gray}
  For stakeholders\par}
\vspace{1.0cm}
\color{rulecolor}\rule{0.8\textwidth}{1.2pt}
\vfill
{\normalsize\color{gray}\today}
\end{titlepage}

% ─────────────────────────────────────────────────────────────────────────────
\section{Protocol Overview}

This protocol allows $N$ parties to simultaneously exchange assets across
different blockchains, without any trusted intermediary.
The parties form a directed cycle: each party deposits funds that the next
party in the cycle can claim, and in turn claims the funds deposited by the
previous party.

\medskip\noindent\textbf{Example.}\;
Three parties $A$, $B$, $C$ want to trade:
$A$ sends Bitcoin to $B$, $B$ sends Cardano ADA to $C$, $C$ sends Bitcoin
to $A$.
The protocol guarantees that under certain assumptions (discussed below) 
either \emph{all three transfers happen} or
\emph{none of them do}: no party can lose funds unless they violate their
own obligations.

\subsection*{How It Works}

\begin{enumerate}
  \item \textbf{Agree.}\;
    All parties agree on the swap terms: who trades what with whom, and the
    safety timing parameters (refund windows).
    One party is designated the \textbf{leader}.

  \item \textbf{Prepare.}\;
    Through a series of message exchanges, the parties jointly construct
    three transactions per transfer:
    a \emph{deposit} (locks funds on-chain),
    a \emph{withdraw} (lets the recipient claim those funds), and
    a \emph{refund} (lets the depositor reclaim if the swap fails).
    The withdraw transaction is incomplete at this stage: it requires a
    shared secret that no single party possesses.

  \item \textbf{Lock.}\;
    Each party broadcasts their deposit transaction.
    Funds are now locked on-chain and can only move via the pre-arranged
    withdraw or refund transactions.

  \item \textbf{Claim.}\;
    The leader collects secret shares from all other parties, completes the
    first withdraw transaction, and broadcasts it on-chain.
    This broadcast publicly reveals the shared secret.
    All other parties extract the secret from the blockchain and complete
    their own withdraw transactions.

  \item \textbf{Refund (if anything goes wrong).}\;
    If the leader never broadcasts the first withdraw, each party's refund
    transaction becomes valid after a pre-agreed block height.
    Every party reclaims their own deposit.
    Refund windows are staggered so that no party can simultaneously
    collect a withdraw and a refund.
\end{enumerate}

The key mechanism is that the act of claiming funds on one chain
\emph{automatically reveals} the information needed for all other parties
to claim on their respective chains.
This is achieved through adaptor signatures, a cryptographic technique that
ties the validity of a signature to the disclosure of a secret.

% ─────────────────────────────────────────────────────────────────────────────
\section{Assumptions and Their Implications}
\label{sec:assumptions}

The protocol's safety guarantees hold only when certain assumptions are met.
This section lists each assumption, explains why it is needed, and
describes the consequences when it is violated.

\subsection{Cyclic Swap Topology}

\textbf{Assumption.}\;
The parties must form a single directed cycle: each party deposits into
exactly one transfer and claims from exactly one transfer.

\textbf{Why.}\;
The staggered refund-window mechanism that prevents double-spending
(collecting both a withdraw and a refund) relies on a strict linear
ordering of the parties.
In non-cyclic topologies (e.g.\ a party that deposits into two transfers),
this mechanism wouldn't work.

\textbf{Implication.}\;
The protocol cannot express arbitrary multi-party trades.
%Any trade must be decomposable into a single cycle.

\subsection{Timely Chain Access}

\textbf{Assumption.}\;
Every party must have \textbf{timely chain access} throughout the protocol:
the ability to both \emph{read} the current state of every relevant
blockchain and \emph{write} transactions that are included in a block,
all within a bounded number of blocks.

Specifically, after the trigger transaction is broadcast, each non-leader
party must be able to (1)~observe the trigger at the required finality
depth, (2)~construct and sign their own withdraw transaction, and
(3)~have that transaction confirmed --- all within $\Delta$ blocks (the
\textbf{claim propagation time}).

\textbf{Why.}\;
The protocol does not tolerate situations where a party is unable to
interact with the chain in time.
Both consistency (all parties see the same chain state) and availability
(all parties can get transactions included) are required.

\textbf{What can break timely chain access:}
\begin{itemize}
  \item \textbf{Party downtime} --- node crash, key unavailability,
    operator error.

  \item \textbf{Network-level censorship} --- miners or validators
    selectively exclude a party's transactions from blocks.
    On PoW chains a mining cartel can censor; on PoS chains a colluding
    validator set can achieve the same effect.

  \item \textbf{Chain congestion} --- sustained high demand delays
    transaction inclusion beyond $\Delta$ blocks even without targeted
    censorship.
    This is economically equivalent to censorship.

  \item \textbf{Network partitions} --- a party is unable to communicate
    with blockchain nodes despite being locally operational.

  \item \textbf{Chain reorganisation} --- a confirmed transaction is
    reversed, resetting the party's view of finality and potentially
    consuming $\Delta$ blocks before re-confirmation.
\end{itemize}

\textbf{Implication.}\;
A party that loses chain access after the Lock phase risks losing its
deposit: another party may claim the locked funds via the withdraw path,
while the affected party cannot claim its own output.
In practice, parties should use monitoring infrastructure (watchtowers,
redundant nodes) and $\Delta$ must be set conservatively for worst-case
inclusion delays.
The protocol offers no defence against a censoring majority of block
producers collaborating with a malicious participant.

\subsection{Sufficient Refund Window Gaps}

\textbf{Assumption.}\;
Adjacent refund windows must be separated by at least $\Delta$ blocks:
$W_i - W_{i+1} \geq \Delta$ for all $i$.

\textbf{Why.}\;
If the gap is too small, a malicious leader can time the trigger broadcast
so that a party's refund window opens before it has had enough time to
confirm its withdraw transaction.
The victim refunds their own deposit (believing the swap failed), but the
leader also claims the victim's locked funds via the withdraw path.
The victim loses their deposit and receives nothing.
Any of the threats to timely chain access listed in the previous section
can cause the effective propagation time to exceed $\Delta$.

\textbf{Implication.}\;
The total lock-up time grows linearly with $N$: the leader's refund window
must be at least $(N{-}1) \times \Delta$ blocks after the earliest window.
Choosing $\Delta$ is a trade-off between safety and capital efficiency;
\textbf{underestimating $\Delta$ is a safety-critical error.}
$\Delta$ should be computed per chain-pair (using the worst case) and
agreed upon by all parties in the Config phase.

\subsection{Chain Finality}

\textbf{Assumption.}\;
The protocol treats ``transaction confirmed'' as a binary oracle.
The definition of finality (e.g.\ number of confirmations) is
chain-specific and external to the protocol.

\textbf{Implication.}\;
A chain reorganisation that reverses a confirmed deposit after the claim
phase has begun can break atomicity irrecoverably.
Cross-chain swaps are only as safe as the weakest finality guarantee among
the chains involved.

\subsection{Schnorr-Compatible Signatures}

\textbf{Assumption.}\;
All participating blockchains must support Schnorr-compatible signature
verification (or an equivalent via smart contracts).

\textbf{Why.}\;
The protocol relies on adaptor signatures, which require the algebraic
structure of Schnorr signatures (specifically, the linearity of nonce
randomisation).
ECDSA-based adaptor signatures exist but have different security properties
and are less mature.

\textbf{Implication.}\;
Bitcoin supports Schnorr natively since Taproot (November 2021).
Cardano uses EdDSA natively, which is not directly compatible; a
smart-contract bridge is required to verify Schnorr signatures on Cardano.
Chains without Schnorr support or smart-contract capability cannot
participate.

\subsection{Cryptographic Soundness}

\textbf{Assumption.}\;
The adaptor signature scheme, the MuSig2 multi-signature scheme, and their
composition are assumed correct and secure.

\textbf{Why.}\;
If the underlying cryptography is broken (e.g.\ an adversary can forge
adaptor signatures or extract the aggregate secret without the trigger),
the protocol's guarantees are void.

\textbf{Implication.}\;
MuSig2 has a published security proof, but the composition of MuSig2 with
adaptor signatures in the specific way this protocol uses them has not been
independently verified in a peer-reviewed setting.
This is a residual cryptographic risk.

\subsection{Mutual Exclusion of Withdraw and Refund}

\textbf{Assumption.}\;
For each transfer, at most one of the withdraw and refund transactions can
be confirmed.
Once one succeeds, the other is permanently invalid.

\textbf{Why.}\;
Without this guarantee, a party could collect both: claim the withdraw
(receiving the counterparty's funds) and then also claim the refund
(recovering their own deposit), stealing funds outright.

\textbf{Implication.}\;
On UTXO-based chains (Bitcoin, Cardano) this holds by construction:
spending a UTXO consumes it.
On account-based chains (e.g.\ Ethereum), a smart contract must
explicitly enforce this exclusion.
Any chain integration must be audited for this property.

\subsection{Griefing and Optionality}
\label{sec:griefing}

\textbf{Assumption.}\;
The protocol guarantees that no honest party loses principal, but it does
not prevent an adversary from forcing all parties to wait out their refund
windows, nor does it prevent the leader from exploiting the temporal gap
between commitment and completion.

\textbf{Why --- griefing.}\;
Any party can abort after the Lock phase by simply withholding their
adaptor secret.
All parties eventually recover their deposits via refund, but their funds
are locked until each party's refund window opens.

\textbf{Why --- optionality.}\;
Between the Lock phase (all funds committed on-chain) and the leader's
decision to broadcast the trigger, time passes and asset prices move.
The leader can observe these movements and choose to complete the swap (if
prices moved in their favour) or let it expire via refund (if prices moved
against them).
The non-leader parties are locked in and cannot exit.
Han et al.~\cite{han2019optionality} prove that this is structurally
equivalent to a premium-free American call option granted to the leader at
the non-leaders' expense, and estimate the fair option premium at 2--3\%
of the swap value for typical cryptocurrency volatility.
This is a structural property of all
atomic swap protocols in which one party
commits last.

\textbf{Implication.}\;
Griefing imposes asymmetric costs: a party with an early refund window can
force the leader to wait at low cost to themselves.
Optionality transfers expected value from non-leaders to the leader, since
the leader can choose to complete or abort based on price movements.
The protocol has no built-in penalty for griefing and no premium mechanism
for optionality; mitigation (reputation systems, bonds, trigger deadlines,
option premiums) is outside the protocol scope.

\subsection{No Upper Bound on Transaction Validity}

\textbf{Assumption.}\;
Bitcoin (and most blockchains) only support a \emph{lower} bound on when a
transaction becomes valid (``not before block $n$''), not an upper bound
(``not after block $m$'').

\textbf{Why.}\;
This means that once a withdraw transaction is valid, it remains valid
\emph{forever}.
There is no way to make it expire.

\textbf{Implication.}\;
Once the refund window opens, both withdraw and refund are simultaneously
valid (mutual exclusion ensures only one succeeds).
A party must broadcast their refund promptly; delay risks the counterparty
completing the withdraw first.
This reinforces the timely chain access requirement (Section~2.2).

\subsection{Pre-Determined Leader}

\textbf{Assumption.}\;
One party is designated as the leader before the protocol begins.
The leader triggers the claim phase by broadcasting the first withdraw
transaction.

\textbf{Why.}\;
The refund windows are anchored to the leader's position: the leader's
refund window opens last, giving all other parties the maximum time to
observe the trigger and claim.
Changing the leader mid-protocol would invalidate the window ordering.

\textbf{Implication.}\;
Leader selection must happen out-of-band or via a pre-protocol election
mechanism (not specified).
The leader role carries both costs and benefits:
\begin{itemize}
  \item \textbf{Cost:}\;
    The leader's funds are locked for the longest duration, and the leader
    bears the highest opportunity cost if the swap is aborted (see
    Section~\ref{sec:griefing}).
  \item \textbf{Benefit:}\;
    The leader holds the free option described in
    Section~\ref{sec:griefing}: they can observe price movements after all
    funds are locked and choose whether to complete or abort.
    This optionality transfers expected value from non-leaders to the
    leader.
\end{itemize}
These two effects pull in opposite directions: griefing risk discourages
taking the leader role, while optionality rewards it.
Whether the net incentive is positive or negative depends on asset
volatility, lock-up duration, and the number of parties.
The protocol does not balance these incentives; rational leader selection
is an open design problem.

% Note that this is a protocol-design choice rather than a fundamental
% limitation: a leader election mechanism
% could be layered on top without changing the core protocol.

\bibliographystyle{plain}
\bibliography{references}

\end{document}
